% page 568 du fac-simile (image p-567 du scan)
% bloc 160.0 x 250.6 mm ; marge gauche 8.0 mm
\pgc{-12.700mm}{0.000mm}{
\l{}
\l{}
\l{}
\l{}
\l{\cel{3}560.}
\l{}
\l{\sou{sucedar}. (trans.) Venar sen interrupteso, sen diskontinueso,}
\l{\cel{3}pos (dop) ulu, ulo. - DEFIS.}
\l{}
\l{\sou{sucesar}. (\sou{netrans., pri}). Abutar fortunoze.- EFIS.}
\l{}
\l{\sou{suciar}. (\sou{netrans., pri)}.\cel{2}Pensar tre ofte e kelke-desquiete pri}
\l{\cel{3}la bona stando di (ulu, ulo). - F.}
\l{}
\l{\sou{sucino}. Speco di rezino fosila, flava, diafana, odoroza, qua}
\l{\cel{3}aquiras, per froto, ula karakterizivi elektrala. - FI.}
\l{}
\l{\sou{sudo}. Un ek la quar punti kardinala, ta qua opozesas a nordo.}
\l{\cel{3}DFIRS.}
\l{}
\l{\sou{sudario}. (\sou{epoki antiqua, en Roma}). Tuko quan la Romani uzis}
\l{\cel{3}por vishar sua vizajo sudorifanta. - L.}
\l{}
\l{\sou{sudoro}. Produkturo de la transpiro di la pelo, kondensita aden}
\l{\cel{3}guteti. - EFIS.}
\l{}
\l{\sou{suficar}. (\sou{netrans., a, pri}). Furnisar sat multe (por satis-}
\l{\cel{3}facar la bezoni, por saciar la deziri.) - EFIS.}
\l{}
\l{\sou{sufixo}. (gram.)\cel{2}Afixo qua, en vorto, adjuntesas dop o pos la}
\l{\cel{3}radiko, por determinar la naturo o senco di olca. - DEFIRS.}
\l{}
\l{\sou{suflar}. I. (\sou{netrans}.) (a) Pulsar, sendar, aero per sua boko,}
\l{\cel{3}(b) Igar ekirar de sua boko, per expiro, la aero respirita,}
\l{\cel{3}- II. (\sou{trans.})\cel{2}(a) Pulsar, sendar, aero adsur (ulo) ;}
\l{\cel{3}(b) Pulsar, sendar,aero aden (ulo). - III. (\sou{trans., ad ulu})}
\l{\cel{3}(\sou{Teatro}) Dicar (ad ulu), tre deslaute, susure, kande lu}
\l{\cel{3}falias memorar lo, to quon lu devas repetar laute. - DFI.}
\l{}
\l{\sou{sufleo}. (\sou{koquarto}). Mikra kuko ek pasto lejera, di qua la in-}
\l{\cel{3}ternajo esas vakua, koquita en mulduro, e qua manjesas}
\l{\cel{3}kom entremeso. - F.}
\l{}
\l{\sou{sufokar}. I. (\sou{trans.)}\cel{2}Haltigar (dum tempo kelketa) la respiro}
\l{\cel{3}(di ulu) per surprizo, dolorigo, iracigo, e c.). - II.}
\l{\cel{3}\sou{(netrans.)}\cel{2}Koaktesar a cesar sua respiro (dum tempo kelk-}
\l{\cel{3}eta) pro surprizeso, doloro, iraco, ridego, singluti, e c.)}
\l{\cel{3}- EFIS.}
\l{}
\l{\sou{sufrar}. (\sou{netrans., de, pro)}.\cel{2}Tolerar sua doloro (fiziolo-}
\l{\cel{3}giala o psikala) dum tempo plu o min multa. EFIS.}
\l{}
\l{\sou{sufragano}. (\sou{eklezio katol.}) Episkopo subordinita ad arkiepis-}
\l{\cel{3}kopo. - DEFIRS.}
\l{}
\l{\sou{sufrajio}. (\sou{liturgio katol.}) I. Prego, da eklezio, da la santi,}
\l{\cel{3}por la religiani. - II. Prego, ye la fino di la laudi e di}
\l{\cel{3}la vespri, por obtenar la interceso da la santi. -}
}
